\chapter{Introduction}
\label{chap:introduction-en}

This chapter situates the difficulties involved in sustaining the facilitation of asynchronous academic discussions when participation grows or the instructor must attend to several threads during the same period. It then sets out the objectives and research questions, describes the work plan and workflow, and closes with the document structure.

\section{Context}

According to social constructivism, learning happens through interaction, in which students construct knowledge by discussing meaning with other students \parencite{WooReeves2007}. In higher education, debates are one of the techniques that best exploit this dynamic because they increase critical thinking and collaborative learning \parencite{Brown2015}, give quieter students a voice, and create a more equitable space for participation than a purely lecture-based class \parencite{BineyEkeyi2018,WooReeves2007}. These discussions may take place in the classroom, where an instructor guides the exchange in real time, or outside it through digital platforms. Outside the classroom, the work is usually asynchronous, so participants contribute at different times without needing to be in the same place.

The digital platforms that support this asynchronous work have grown significantly in Europe in recent years. In 2021, 100\% of institutions in the European Higher Education Area already used some form of digital learning \parencite{GaebelEtAl2021}. A survey published that year found that blended learning was used in 75\% of institutions and that 36\% offered at least one fully online degree program \parencite{GaebelEtAl2021}. These platforms are known as learning management systems (LMSs) and, depending on the institution, may provide the main course environment or complement in-person education. Moodle, Canvas, Open edX, and Coursera, for example, allow content, activities, and discussion forums to be organized in a single environment, with instructor-led or self-paced courses.

The 2020 pandemic accelerated the adoption of digital learning. During the first four months of the crisis, several institutions reported more progress in digitalization than in the four previous years \parencite{GaebelEtAl2021}. Consequently, interaction between students came to happen almost exclusively through digital tools, and discussion forums became one of its main mechanisms in online education \parencite{Rovai2007,MasseyEtAl2019}.

In this context, discussion forums offer forms of participation that are difficult to sustain in an in-person classroom, including scale, turn-taking equity, and scheduling flexibility. A study of 5,104 students in parallel discussions through Moodle illustrates the volume of participants that can be brought together in this type of environment \parencite{ZulfikarEtAl2019}. They also give each student more opportunities to participate than the in-person format, where the instructor can take approximately 65\% of talk time \parencite{SullivanPratt1996}, since everyone can contribute regardless of location or schedule availability \parencite{MasseyEtAl2019}. Forums also support forms of assessment based on the quality of contributions, which can measure critical thinking more directly than a traditional exam \parencite{HoSwan2007,Rovai2007}.

These possibilities have extended the use of forums across settings of very different scales, including massive open online courses (MOOCs). As the number of participants and active discussion threads grows, so does the work required to facilitate the discussions. This difficulty provides the starting point for the next section.

\section{Motivation}

The facilitation of asynchronous discussions falls mainly on the instructor, who performs organizational, intellectual, and social functions, according to Paulsen's classification as reported by Abdous \parencite{Abdous2011}. In one study, 79.3\% of contributions reached deep critical thinking \parencite{LimEtAl2011}. In another, structured roles increased knowledge construction in the groups studied, particularly at the beginning of the discussions \parencite{DeWeverEtAl2010}. Rovai proposes 20--30 students as a practical estimate for an active discussion forum and warns that a larger number may overburden facilitation \parencite{Rovai2007}. In postgraduate courses, 57\% of the recorded instructor-presence behaviors corresponded to social presence \parencite{RichardsonEtAl2015}; corrective feedback represented only 4.9\% of posts in another study \parencite{BlignautTrollip2003}; and facilitators reported difficulty preserving depth across concurrent discussions \parencite{Mansour2024}. The problem becomes more acute in MOOCs, where a single course may include thousands of students \parencite{GuEtAl2021,YuEtAl2024maic}.

Available reviews document that these applications concentrate mainly on personalized tutoring, feedback generation, and curriculum design assistance \parencite{ChuEtAl2025}. However, the facilitation of group discussions remains less explored. Existing natural language processing (NLP) work concentrates on classification tasks, such as toxicity detection; the production of interventions intended to improve discussion quality receives less attention \parencite{KorreEtAl2025}. These reviews more often describe uses aimed at structuring learning conditions than active interventions in group dynamics \parencite{CasebourneEtAl2025}. Some work addresses discussion facilitation directly, including a case-based reasoning system for generating suggestions in online forums \parencite{GuEtAl2021}, and MAIC (\textit{Massive AI-empowered Course}), which organizes teaching functions in massive environments using multiple LLM agents \parencite{YuEtAl2024maic}.

Facilitation requires deciding when a discussion needs support, what type of intervention may be appropriate, and how to formulate it. This work addresses that need with a system that analyzes asynchronous academic discussion threads and proposes interventions based on facilitation criteria.

\section{Objectives}

The objective is to design and implement a system based on LLMs and context engineering techniques to analyze and facilitate asynchronous, text-based academic discussions from different sources. The system uses thread content and criteria derived from the pedagogical literature to decide whether to intervene and to produce a contextualized response.

The specific objectives are the following:

\begin{itemize}
  \item Design and implement an LLM-based agent system, organized in stages, capable of producing interventions coherent with the discussion topic and with conversational continuity, within reasonable time and resource consumption.
  \item Define and implement a context engineering strategy grounded in the literature that organizes thread content, facilitation roles, and techniques, using performance criteria derived from how an instructor effectively facilitates a discussion.
  \item Define a common input contract and use it with reproducible scenarios, six anonymized historical threads preserved in the repository, and a proof of concept with Open edX as a source.
  \item Evaluate the system's behavior in scenarios derived from the literature and in imported threads, examining the classifications produced and the estimated adequacy of the generated interventions.
\end{itemize}

\section{Work plan}

The work partially follows the \textit{Design Science Research} process of \textcite{PeffersEtAl2007}.

The project uses six milestones (table~\ref{tab:milestones-en}) covering problem exploration, design and implementation, inputs and outputs, technical evaluation, and deployment.

\begin{table}[h]
\centering
\small
\renewcommand{\arraystretch}{1.15}
\begin{tabularx}{\textwidth}{@{}Z{3.8cm}Y@{}}
\toprule
\textbf{Milestone} & \textbf{Description} \\
\midrule
1. Problem definition & Formalization of the problem, objectives, requirements, and methodological approach. \\
2. Research and exploration & Proof-of-concept tests to reduce technical uncertainty before design. \\
3. Implementation & Design and implementation of the system based on the findings from the previous milestone. \\
4. Inputs and outputs & Common input contract, execution result, and proof of concept with Open edX. \\
5. Quality assurance and feedback & Evaluation, result collection, and preparation of the final documentation. \\
6. Deployment and infrastructure & Operation of the environments used for the experiments and technical demonstration. \\
\bottomrule
\end{tabularx}
\caption{Project milestones and their relationship to the research, implementation, and evaluation phases.}
\label{tab:milestones-en}
\end{table}

Figure~\ref{fig:roadmap-github-project} provides the corresponding timeline recorded in GitHub Project and shows where the research, implementation, integration, and evaluation tasks overlapped.

\section{Research questions}

The research began with exploratory questions before any design decisions were made. The initial hypothesis was that asynchronous discussions on educational platforms might not produce the same type of learning as an in-person discussion. It was not yet clear what could explain this difference or how it might be addressed.

\subsection{Exploratory phase}

The first questions aimed to understand that dynamic:

\begin{itemize}
  \item What happens in an in-person classroom that does not happen in an asynchronous discussion? Why doesn't the social dynamic of learning carry over to text forums?
  \item What makes a facilitator effective? What concrete actions does an instructor perform when intervening in a discussion?
  \item What does the intervention process look like from the inside? What information does a facilitator need to make a decision?
  \item Can the intervention process be described through observable and controllable decisions?
\end{itemize}

These questions guided the literature review and the initial technical exploration. The fourth question had a particular influence on the design. Evidence from intelligent tutoring systems \parencite{VanLehn2011} and the literature on facilitation roles \parencite{PilkingtonWalker2003} support framing the intervention process as a sequence of steps with recordable intermediate decisions.

The state-of-the-art review in Chapter~\ref{chap:estado-arte} examines the instructor functions and facilitation roles described in the literature. On this basis, the design questions focused on how to represent those functions in the system.

\subsection{Design phase}

Once the problem was delimited, the questions became more concrete:

\begin{itemize}
  \item How can facilitation roles be operationalized as LLM agents in a discussion analysis sequence?
  \item What criteria allow automatic detection of a discussion's state to determine whether and when to intervene?
  \item What criteria allow evaluation of whether an AI-generated intervention is pedagogically appropriate?
  \item Can the system receive discussions from an educational platform without coupling the sequence to its data model?
\end{itemize}

These questions structure the development chapter, since each one corresponds to a system component or a documented design decision.

\section{Workflow}

The project uses an R\&D workflow inspired by agile practices, with kanban as the tracking tool but without fixed-length iterations. Some questions are resolved in days, whereas others require weeks of exploration before a decision can be made.

In practice, the work is organized in discovery cycles, so the findings from each phase inform the next. Tasks, the GitHub board, and Architecture Decision Records (ADRs), available in the project repository\footnote{\href{https://github.com/magrimal/tfm-2026-discussion-moderation}{Project repository}; \href{https://github.com/users/magrimal/projects/3/views/1}{work board}}, document the exploration and design decisions. Milestones delimit the phases, and each phase ends when the questions that define it have been answered.

Technical uncertainties are examined through exploratory tests or proof-of-concept implementations before design decisions are consolidated. These tests have a limited scope and guide subsequent development.

\subsection{Use of AI assistants}

The implementation was carried out with support from LLM-based programming assistants, in particular Claude Code and Codex. They were used to draft code and documentation, explore alternatives, locate inconsistencies, and propose tests. NotebookLM and Google Scholar Labs were used to explore terminology and locate literature. Grammarly and ChatGPT supported grammatical review and iterative improvement of drafts. The work followed cycles of proposal, review, testing, and correction.

The problem was defined and proposals were accepted or rejected; responsibility for decisions, experiments, and the final text was retained throughout. No percentage of human- versus AI-authored code was calculated; Git history documents the implementation process in enough detail to reconstruct it.

\section{Document structure}

This document is organized into seven chapters, including the required English versions of the introduction and conclusions. Chapter~\ref{chap:introduccion} presents the problem, objectives, research questions, and work plan in Spanish, and this chapter provides its English version. Chapter~\ref{chap:estado-arte} reviews discussion facilitation and related educational systems. Chapter~\ref{chap:desarrollo} describes the architecture, facilitation model, analysis sequence, inputs and outputs, context strategy, inspection panel, and deployment environments. Chapter~\ref{chap:resultados} reports the experiments. Finally, Chapter~\ref{chap:conclusiones} presents the conclusions, limitations, and future work in Spanish, and Chapter~\ref{chap:conclusions-en} provides the corresponding English version. Appendix~\ref{app:flujo-prototipo} preserves the complete interface usage sequence.
